\clearpage\begingroup\SourceLeaf{425b}\ArithmeticHead{94}{L}{STÄCKEL, GAUSS ALS GEOMETER.}
\renewcommand{\thefootnote}{\arabic{footnote})}\setcounter{footnote}{1}
\makeatletter\tagsleft@true\let\veqno\leqno\makeatother
\begin{center}
\textbf{30.}\par\smallskip
\textls[100]{Vorarbeiten zu den Allgemeinen Untersuchungen}\par
\textls[100]{über die krummen Flächen} (1822—1825).
\end{center}
\noindent\hspace*{\parindent}\textsc{Gauss} hatte der Kopenhagener Preisschrift als Kennwort den Ausspruch
\textsc{Newtons} mitgegeben: Ab his via sternitur ad maiora; diese Worte bilden
den Schluss der 1704 als Anhang zur Optik veröffentlichten Abhandlung \textit{De
quadratura curvarum}, in der \textsc{Newton} ältere Untersuchungen bekannt gab, die
ihn zur Fluxionsrechnung geführt hatten\footnote{\textsc{J. Newton}, \textit{Opuscula mathematica}, rec. I. Castillioneus, vol. 1, Lausanne und Genf 1744, S. 244; es heisst wörtlich: »Et his principiis via ad maiora sternitur«.}.
\clearpage\SourceLeaf{426}\ArithmeticHead{95}{R}{DIE ALLGEMEINE LEHRE VON DEN KRUMMEN FLÄCHEN.}\setcounter{footnote}{0}
\indent Was waren die grösseren Dinge, zu denen die konforme Abbildung den
Weg bahnte? Eine Andeutung findet man im Art. 4 der Preisschrift: »Wenn
überdies [das Vergrösserungsverhältnis bei der Abbildung] \(m=1\) ist, wird eine
vollkommene Gleichheit [der einander entsprechenden Linienelemente] stattfinden,
und die eine Fläche sich auf die andere abwickeln lassen« (W. IV,
S. 195). Dass die Lehre von der Abwicklung oder Biegung der krummen
Flächen gemeint war, beweist eine Aufzeichnung, die \textsc{Gauss} am 13. Dezember
1822, zwei Tage, nachdem er die Beantwortung der Preisfrage an \textsc{Schumacher}
abgesandt hatte, begonnen und am 15. Dezember beendet hat (W. VIII,
S. 374—384). Sie führt den Titel: Stand meiner Untersuchung über die
Umformung der Flächen und zeigt, dass er damals für den besonderen Fall,
wo das Quadrat des Linienelementes vermöge der konformen Abbildung auf
die Form
\begin{equation}ds^2=m(dp^2+dq^2)\tag{1}\end{equation}
\noindent gebracht ist, die Rechnungen durchgeführt hat, die sich für den allgemeinen
Fall, wo
\begin{equation}ds^2=E\,dp^2+2F\,dp\,dq+G\,dq^2\tag{2}\end{equation}
\noindent ist, in den Artt. 9 und 10 der \textit{Disq. gen.} finden. Das Endergebnis besteht in
dem Lehrsatz, dass das Krümmungsmass der Fläche allein durch die Funktion
\(m(p,q)\) und deren erste und zweite partielle Ableitungen nach \(p\) und \(q\) ausgedrückt
werden kann. Hieraus folgt sogleich, dass das Krümmungsmass bei
den Biegungen einer Fläche erhalten bleibt.
\par
Wir haben gesehen, dass \textsc{Gauss} das »schöne Theorem« von der Erhaltung
des Krümmungsmasses oder genauer von der Erhaltung der Gesamtkrümmung
solcher Flächenstücke, die durch Biegung aus einander hervorgehen, bereits
im Jahre 1816 besass. Wenn man annimmt, dass er den vorstehenden aus
der \textls[100]{konformen Abbildung} fliessenden Beweis, der in der Aufzeichnung
vom Dezember 1822 als Ziel der Untersuchung erscheint, in der Zeit zwischen
1816 und 1822 gefunden hat, so entsteht die Frage, welches die ursprüngliche
Quelle für das Theorem gewesen ist. Aufzeichnungen aus der Zeit vor 1816,
die sich darauf beziehen, sind nicht vorhanden, es lässt sich jedoch sehr wahrscheinlich
machen, dass die Lehre von den \textls[100]{kürzesten Linien} auf krummen
Flächen den Zugang eröffnet hat.
\clearpage\SourceLeaf{427}\ArithmeticHead{96}{L}{STÄCKEL, GAUSS ALS GEOMETER.}\setcounter{footnote}{0}
\indent Die stärksten Gründe für diese Behauptung ergeben sich aus einem ersten
Entwurfe der \textit{Disq. gen.}, der den Titel führt: »\textit{Neue allgemeine Untersuchungen
über die krummen Flächen}« und der aus den letzten Monaten des Jahres 1825
stammt (W. VIII, S. 408—442). Es empfiehlt sich daher, zunächst die Entstehung
dieses Entwurfs zu schildern und jene Frage im Zusammenhang mit
dem Bericht über dessen Inhalt zu erörtern.
\par
Schon am 28. Juli 1823 hatte \textsc{Gauss}, an die kürzlich erfolgte Erteilung
des Kopenhagener Preises anknüpfend, zu \textsc{Olbers} bemerkt: »Sollte ich in
diesem Leben noch einmal in eine den Arbeiten günstigere Lage kommen,
so werde ich diese Abhandlung [die Preisschrift] mit als Teil einer viel ausgedehnteren
Untersuchung verarbeiten« (Br. G.-O., 2, S. 252). Er meinte
damit ein grösseres, die Theorie und die Praxis der höheren Geodäsie behandelndes
Werk. Ein solcher Plan wird ausdrücklich in dem Brief an \textsc{Olbers}
vom 9. Oktober 1825 erwähnt. »Ich habe dieser Tage angefangen, in
Beziehung auf mein künftiges Werk über Höhere Geodäsie einen (sehr) kleinen
Teil dessen, was die krummen Flächen betrifft, in Gedanken etwas zu ordnen.
Allein ich überzeuge mich, dass ich bei der Eigentümlichkeit meiner ganzen
Behandlung des Zusammenhanges wegen gezwungen bin, \textls[100]{sehr weit auszuholen},
sodass ich sogar meine Ansicht über die Krümmungshalbmesser
bei \textls[100]{planen} Kurven vorausschicken muss. Ich bin darüber fast zweifelhaft
geworden, ob es nicht geratener sein wird, einen Teil dieser Lehren, der ganz
rein geometrisch (in analytischer Form) ist und Neues mit Bekanntem gemischt
in neuer Form enthält, erst besonders auszuarbeiten, es vielleicht von dem
Werke abzutrennen und als eine oder zwei Abhandlungen in unsere Commentationes
einzurücken. Indessen kann ich noch vorerst die Form der Bekanntmachung
auf sich beruhen lassen und werde einstweilen in dem zu Papier
bringen fortfahren« (W. VIII, S. 397, IX, S. 376)\footnote{Vgl. auch den Brief an \textsc{Pfaff} vom 21. März 1825: »Nach Beendigung der Messungen werde ich darüber ein eigenes Werk, vermutlich von bedeutender Ausdehnung, ausarbeiten« (W. X 1, S. 250).}.
\par
Die Briefe an \textsc{Schumacher} vom 21. November 1825 (W. VIII, S. 400) und
an \textsc{Hansen} vom 11. Dezember 1825 (Brief im \textsc{Gauss}-Archiv) zeigen, dass \textsc{Gauss}
bis gegen Ende des Jahres an dem Entwurf gearbeitet hat. Mit der Darstellung
des Krümmungsmasses bei geodätischen Polarkoordinaten, für die
\begin{equation}ds^2=dp^2+G\,dq^2\tag{3}\end{equation}
\clearpage\SourceLeaf{428}\ArithmeticHead{97}{R}{DIE ALLGEMEINE LEHRE VON DEN KRUMMEN FLÄCHEN.}\setcounter{footnote}{0}
\noindent wird, hat er abgebrochen, augenscheinlich, weil er jetzt erkannte, dass es
möglich sei, eine entsprechende Formel für die allgemeine Form (2) des Quadrats
des Linienelementes aufzustellen. Bei der wirklichen Durchführung
dieses Gedankens, an die er sogleich ging, ist \textsc{Gauss} auf grosse Schwierigkeiten
gestossen. Er hat sie erst im Herbst 1826 überwunden. Hierüber
wird in der nächsten Nummer berichtet werden, in dieser Nummer wenden
wir uns zu den \textit{Neuen allgemeinen Untersuchungen}.
\par
Wie \textsc{Gauss} in dem Brief an \textsc{Olbers} vom 9. Oktober 1825 angekündigt
hatte, beginnen die \textit{Neuen allgemeinen Untersuchungen über die krummen Flächen}
mit den ihm eigentümlichen Ansichten über die Krümmung ebener Kurven
(artt. 1—6). Zwei Punkte sind dabei wesentlich, erstens dass er gerichtete
Gerade einführt und so die Frage der Vorzeichen klärt, zweitens, dass die
Krümmung der Kurven mittels derjenigen Abbildung auf den Kreis vom
Halbmesser Eins eingeführt wird, bei der Punkte mit parallelen Normalen
einander entsprechen.
\par
Zum Raume übergehend bringt \textsc{Gauss} zunächst (artt. 7—8) sieben einleitende
Sätze, die später in die artt. 1, 2 und 4 der \textit{Disq. gen.} aufgenommen
worden sind; sie dienen dazu, die Abbildung der krummen Fläche auf die
Kugel vom Halbmesser Eins mittels paralleler Normalen vorzubereiten. Das
vorletzte Theorem ist neu; es findet sich auch im Handbuch 19 Be, S. 78 und
stammt aus der Zeit um 1810.
\par
Es folgt (artt. 9—11) die Untersuchung des Verhaltens einer krummen
Fläche in der Umgebung eines regulären Punktes. \textsc{Gauss} benutzt hier nicht
wie in den \textit{Disq. gen.} (art. 8) das Verfahren der Reihenentwicklung, sondern
betrachtet die Schnittkurven der Fläche mit dem Büschel der durch den betrachteten
Punkt gehenden Ebenen; er gelangt daher hier auch zu dem Satze
von \textsc{Meusnier}, der in den \textit{Disq. gen.} nicht vorkommt.
\par
Am Schluss des art. 11 wird die Abbildung der krummen Fläche auf die
Einheitskugel mittels paralleler Normalen gelehrt. Von hier aus gelangt man,
in Verallgemeinerung der bei den ebenen Kurven angestellten Überlegungen,
zu den Begriffen der Gesamtkrümmung eines Flächenstückes und des Krümmungsmasses,
das einem Flächenpunkte zugeordnet ist. Die einem Flächenstück
entsprechende Area auf der Einheitskugel wird hier noch nicht als
deren Gesamtkrümmung bezeichnet. Dieser Name ist also wohl erst später
\par\smallskip\noindent{\small X\textsubscript{2} Abh. 4.\hfill 13}\par
\clearpage\SourceLeaf{429}\ArithmeticHead{98}{L}{STÄCKEL, GAUSS ALS GEOMETER.}\setcounter{footnote}{0}
\noindent entstanden. In dem Brief an \textsc{Olbers} vom 20. Oktober 1825 sagt \textsc{Gauss},
seine Untersuchungen bezögen sich auf eine Menge von Gegenständen, die
er nicht anführen könne, »weil die Begriffe davon nicht gangbar sind und
selbst noch keine Namen dafür existieren« (Br. G.-O. 2, S. 431, W. VIII,
S. 398). Endlich wird der Zusammenhang zwischen dem Krümmungsmass und
den beiden Hauptkrümmungen entwickelt.
\par
Im Unterschied gegen die \textit{Disq. gen.} wendet sich \textsc{Gauss} nunmehr sogleich
zu den \textls[100]{kürzesten Linien}, die auf der betrachteten krummen Fläche liegen,
und geht hier auch auf ganz andere Art vor als dort.
\par
Die Aufgabe, zwei gegebene Punkte einer krummen Fläche durch die
kürzeste Linie zu verbinden, war 1697 von \textsc{Johann Bernoulli} den Geometern
gestellt worden\footnote{\textsc{Joh. Bernoulli}, Journal des savants, année 1697, S. 394, Opera omnia, Lausanne 1742, t. I, S. 204.}, aber erst 1732 hatte \textsc{Euler} eine Lösung veröffentlicht\footnote{\textsc{L. Euler}, \textit{De linea brevissima in superficie quacunque duo quaelibet puncta iungente}, Comment. acad. sc. Petrop. t. 3 (1728), 1732, S. 110.};
\textsc{Bernoulli} gab sein Verfahren 1742 bekannt\footnote{\textsc{Joh. Bernoulli}, \textit{In superficie quacunque curva ducere lineam inter duo puncta brevissimam}, Opera omnia, Lausanne 1742, t. IV, S. 108.}. Im Laufe des 18. Jahrhunderts
wurden besonders die kürzesten Linien auf dem elliptischen Sphäroide untersucht,
weil sie für die Geodäsie wichtig waren. Diese Kurven wurden daher
als geodätische Linien bezeichnet; erst seit der Mitte des 19. Jahrhunderts ist
es üblich geworden, bei beliebigen krummen Flächen von geodätischen Linien
zu sprechen\footnote{Vgl. \textsc{P. Stäckel}, \textit{Bemerkungen zur Geschichte der geodätischen Linien}, Leipziger Berichte, 1893, S. 444.}.
\par
\textsc{Gauss} war, wie wir gesehen haben (S. 85), schon vor 1816 damit beschäftigt
gewesen, die Lehre von den kürzesten Linien des Sphäroides für die
Zwecke der Geodäsie auszubauen. Er hat sich aber damals auch schon mit
den kürzesten Linien auf beliebigen krummen Flächen beschäftigt, denn in
dem Briefe an \textsc{Schumacher} vom 21. November 1825 schreibt er, seine allgemeinen
Untersuchungen über die krummen Flächen seien durch manchen
glücklichen Fund belohnt worden. »So habe ich zum Beispiel die Generalisierung
des \textsc{Legendreschen} Theorems, dass auf der Kugel die Seiten [eines
kleinen sphärischen Dreiecks] proxime den Sinus der um \(\frac{1}{3}\) des sphärischen
Exzesses verminderten Winkel proportional sind, auf krumme Flächen \textls[100]{jeder}
\clearpage\SourceLeaf{430}\ArithmeticHead{99}{R}{DIE ALLGEMEINE LEHRE VON DEN KRUMMEN FLÄCHEN.}\setcounter{footnote}{0}
\noindent Art (wo die Verteilung \textls[100]{ungleich} geschehen muss), welche ich der Materie
nach schon seit vielen Jahren besessen, aber noch nicht zu möglicher Mitteilung
an andere entwickelt hatte, jetzt auf eine überaus elegante Gestalt
gebracht« (W. VIII, S. 400; vgl. auch den Brief an \textsc{Olbers} vom 20. Oktober
1825, Br. G.-O. 2, S. 431, W. VIII, S. 399). In einer gleichzeitig niedergeschriebenen
Aufzeichnung hat \textsc{Gauss} sein Verfahren angedeutet (W. VIII,
S. 401—405); jene ungleiche Verteilung wird danach bedingt durch die Werte,
die dem Krümmungsmass der Fläche in den Eckpunkten des Dreiecks zukommen.
\par
\textsc{Legendre} hatte sein Theorem von der Zurückführung eines kleinen sphärischen
Dreiecks auf ein ebenes Dreieck mit Seiten derselben Länge 1789
ohne Beweis bekanntgemacht\footnote{\textsc{A. M. Legendre}, \textit{Mémoire sur les opérations trigonométriques dont les résultats dépendent de la figure de la terre}, Histoire de l’Acad., année 1787, Paris 1789, Mémoires, S. 358.} und den Beweis 1798 nachgeholt\footnote{\textsc{A. M. Legendre}, \textit{Résolution des triangles sphériques dont les côtés sont très-petits, pour la détermination d’un arc de méridien}, Note III des Werkes von \textsc{Delambre}, \textit{Méthodes analytiques pour la détermination d’un arc de méridien}, Paris, an VIII; kurz darauf erschien im Journal de l’école polytechnique ein Beweis von \textsc{Lagrange}, \textit{Solutions de quelques problèmes relatifs aux triangles sphériques}, t. II, cah. 6, Paris 1798, S. 270; Oeuvres, t. 7, S. 329. Der Merkwürdigkeit wegen sei hier auf die anmassende Kritik hingewiesen, die \textsc{Kaestner} in seinen \textit{Geometrischen Abhandlungen}, 2. Sammlung, Göttingen 1791, S. 456—458 an dem \textsc{Legendreschen} Theorem geübt hat, vielleicht hat sie zu dem geringschätzigen Urteil beigetragen, das \textsc{Gauss} über \textsc{Kaestner} als Mathematiker gefällt hat.}.
\par
Bei einer Verallgemeinerung auf geodätische Dreiecke beliebiger krummer
Flächen musste der erste Schritt sein, die Winkelsumme eines solchen Dreiecks
zu ermitteln, und nun sehen wir, dass \textsc{Gauss} in den \textit{Neuen allgemeinen
Untersuchungen}, zu denen wir hiermit zurückkehren, nachdem er bewiesen
hat, dass für jeden Punkt einer kürzesten Linie die Schmiegungsebene die
betreffende Flächennormale in sich enthält (art. 12), sogleich zu dem Satze
übergeht, dass die Summe der Winkel eines geodätischen Dreiecks von zwei
Rechten um einen Betrag abweicht, der durch den Inhalt des entsprechenden
Dreiecks auf der Einheitskugel gegeben wird, wenn man deren Oberfläche
gleich acht Rechten setzt.
\par
\textsc{Gauss} schreibt am 21. November 1825, er habe die Generalisierung des
\textsc{Legendreschen} Theorems schon seit vielen Jahren besessen. Man wird daher
annehmen dürfen, dass er die ersten Schritte dazu schon vor 1816 gemacht
hatte, dass er also schon damals den Satz von der Winkelsumme eines geodätischen
\par\smallskip\noindent\hfill13\textsuperscript{*}\par
\clearpage\SourceLeaf{431}\ArithmeticHead{100}{L}{STÄCKEL, GAUSS ALS GEOMETER.}\setcounter{footnote}{0}
\noindent Dreiecks besass, zu dessen Herleitung die Kenntnis der einfachsten
Eigenschaften der kürzesten Linien genügt, sobald man den genialen Gedanken
der Abbildung mittels paralleler Normalen gefasst hat; die Beziehung
der Richtungen im Raume auf die Einheitskugel hat \textsc{Gauss} aber schon im
Jahre 1799 besessen, das zeigt die bereits S. 87 erwähnte Aufzeichnung vom
November 1799. Wir wissen ferner, dass er schon vor 1816 die Biegung
krummer Flächen betrachtet und nach Kennzeichen dafür gefragt hatte, dass
zwei gegebene Flächen durch Biegung aus einander hervorgehen (W. VIII,
S. 372). Bei der Biegung entsteht aber aus einem geodätischen Dreieck
wieder ein geodätisches Dreieck mit denselben Winkeln, es bleibt also die
Winkelsumme erhalten und damit auch die Grösse der Area, die dem geodätischen
Dreieck auf der Einheitskugel bei der Abbildung mittels paralleler
Normalen entspricht. Denkt man sich also ein beliebiges Flächenstück in
geodätische Elementardreiecke (Triangulation) zerlegt, so folgt, dass bei der
Biegung irgend welchen einander entsprechenden Flächenstücken gleich grosse
Flächenstücke auf der Einheitskugel zugeordnet werden, und das ist genau
das »schöne Theorem«. Wird schliesslich, damit man zu einer Funktion des
Ortes auf der Fläche gelangt, in naturgemässer Verallgemeinerung des Begriffes
der Krümmung bei Kurven das Krümmungsmass bei Flächen als der
Grenzwert erklärt, dem das Verhältnis der Area auf der Einheitskugel zu
dem entsprechenden Flächenstück zustrebt, wenn dieses auf den betrachteten
Punkt zusammenschrumpft, so ergibt sich »der wichtige Lehrsatz, dass bei der
Übertragung der Flächen durch Abwicklung das Krümmungsmass an jeder
Stelle unverändert bleibt«, und das ist das Endergebnis der Entwicklungen in
den artt. 13—16 der \textit{Neuen allgemeinen Untersuchungen}. Die hier gegebene
Herleitung wird man mithin als die ursprüngliche, vor 1816 gefundene, dagegen
die Herleitung aus der Form (1) des Quadrates des Linienelements als
die spätere, zwischen 1816 und 1825 entstandene anzusehen haben.
\par
Es folgt der Beweis des Satzes, dass der Ort der Punkte gleicher geodätischer
Entfernung von einem Punkte der Fläche eine Kurve ist, die alle
von dem Punkte ausgehenden geodätischen Linien unter rechtem Winkel
schneidet (art. 17), und den Schluss des Entwurfes bildet der Satz, dass bei
Einführung geodätischer Polarkoordinaten, die dem Quadrate des Linienelements
die Gestalt (3) verleihen, das Krümmungsmass allein durch die
\clearpage\SourceLeaf{432a}\ArithmeticHead{101}{R}{DIE ALLGEMEINE LEHRE VON DEN KRUMMEN FLÄCHEN.}\setcounter{footnote}{0}
\noindent Funktion \(G(p,q)\) und deren erste und zweite partielle Ableitungen nach \(p\)
und \(q\) ausgedrückt werden kann (art. 18).
\par
Damit war ein dritter Beweis für die Erhaltung des Krümmungsmasses
gegenüber Biegungen gefunden. Was aber bei den beiden besonderen Formen
(1) und (3) des Linienelementes gelungen war, musste auch für die allgemeine
Form (2) gelten, das heisst, es musste möglich sein, das Krümmungsmass
allein durch die Funktionen \(E(p,q),\ F(p,q),\ G(p,q)\) und deren erste und
zweite partielle Ableitungen auszudrücken. So lange das nicht geleistet war,
hatte die Lehre vom Krümmungsmass keine befriedigende Gestalt gewonnen,
und daher hat \textsc{Gauss} Ende 1825 den Entwurf beiseite gelegt, »nil fecisse
putans, si quid superesset agendum«.
\clearpage\endgroup
